\section{String-flow vocabulary}

This chapter records the exact vocabulary used in the main paper and in
the Lean repository.

\subsection{Words}

A word is a finite list of natural numbers $t_0,\dots,t_{L-1}$.  The
length is $L$.  The word is valid from $x_0$ when each division
\begin{equation*}
  x_{i+1}=\frac{5x_i+1}{2^{t_i}}
\end{equation*}
is exact.

The Lean names are:
\begin{itemize}[leftmargin=*]
\item \texttt{StringFlow.Word.wordValid};
\item \texttt{StringFlow.Word.wordOrbit};
\item \texttt{StringFlow.wordWeight};
\item \texttt{StringFlow.Word.wordA}.
\end{itemize}

\subsection{Recursive definitions}

The Lean definitions in \texttt{WordWindow.lean} are:
\begin{verbatim}
def wordA : List Nat -> Nat
  | [] => 0
  | t :: ts => 5 ^ ts.length + 2 ^ t * wordA ts

def wordOrbit : List Nat -> Nat -> Nat
  | [], x => x
  | t :: ts, x => wordOrbit ts ((5 * x + 1) / 2 ^ t)

def wordValid : List Nat -> Nat -> Prop
  | [], _ => True
  | t :: ts, x => (5 * x + 1) % 2 ^ t = 0 /\
    wordValid ts ((5 * x + 1) / 2 ^ t)
\end{verbatim}

The paper writes these as $\wordA$, $\wordOrbit$, and the validity
predicate defined in the main text.

\subsection{The valuation lemma}

Every positive integer $n$ decomposes uniquely as
\begin{equation*}
  n=2^{\vtwo{n}}\,\operatorname{oddpart}(n),
\qquad \operatorname{oddpart}(n)\equiv1\pmod2.
\end{equation*}
The exact step weight is the maximal legal weight:
\begin{equation*}
  t=\vtwo{5x+1},
\end{equation*}
while \texttt{wordValid} only requires divisibility by $2^t$.  A legal
word may therefore use a smaller weight, which is why
$\GeneralOrbit$ can reach intermediate states not on the accelerated
orbit.  The Lean names are \texttt{n\_eq\_two\_pow\_mul\_oddPart} and
\texttt{oddPart\_odd\_of\_pos}.

\subsection{Word molecule identity}

For a valid word $w$ from $x_0$,
\begin{equation*}
  2^{\wordWeight(w)}\,\wordOrbit(w,x_0)
    =5^{\operatorname{len}(w)}\,x_0+\wordA(w).
\end{equation*}

\subsection{Proof of the word orbit identity}

Proceed by induction on $w$.  For $w=[]$, both sides equal $x_0$.  For
$w'=t::w$, the inductive hypothesis applied to
$(5x_0+1)/2^t$ gives
\begin{equation*}
  2^{\wordWeight(w)}\,\wordOrbit\!\left(w,\frac{5x_0+1}{2^t}\right)
    =5^{\operatorname{len}(w)}\frac{5x_0+1}{2^t}+\wordA(w).
\end{equation*}
Multiplying by $2^t$ and using $\wordWeight(w')=t+\wordWeight(w)$ and
$\wordA(w')=5^{\operatorname{len}(w)}+2^t\wordA(w)$ gives the identity.

\subsection{Exact orbit bound}

For $n\ge2$,
\begin{equation*}
  \Orbit_7(n)<5^n.
\end{equation*}
If $x<5^k$, then $5x+1<2\cdot5^{k+1}$, so $(5x+1)/2<5^{k+1}$; since
the accelerated step is at most $(5x+1)/2$, induction gives the
bound.  The Lean statement is
\texttt{fiveXPlusOneOrbit\_lt\_five\_pow}.

\subsection{PMI identity}

With $m_j=j\log_2 5-\wordWeight_j$,
\begin{equation*}
  \sum_{j=0}^{L-1}2^{-m_j}
    =\frac{5\wordA(w)}{5^L}.
\end{equation*}
The cleared-denominator form is
\begin{equation*}
  \sum_{j=0}^{L-1}2^{\wordWeight_j}\,5^{L-j}
    =5\wordA(w).
\end{equation*}
For a closed cycle word of period $P$ and cycle state $m$, the Lean
form is
\begin{equation*}
  \mathrm{aTotal5}(P,t)=5m(2^T-5^P),
\qquad T=\wordWeight(w),
\end{equation*}
proved in \texttt{StringFlow.PMI.pmi\_algebraic}.

\subsection{PMI-B}

The PMI-B layer compares bad-prefix counts with the available budget.
For a closed cycle word, the total weight satisfies
\begin{equation*}
  2^{\wordWeight(w)}\,m=5^P\,m+\wordA(w).
\end{equation*}

\subsection{PMI-B derivation}

Write $B$ for the number of bad proper prefixes.  The $j=0$ term of
$\mathrm{aTotal5}$ contributes $5^P$, and each bad prefix contributes
at least another $5^P$, so
\begin{equation*}
  5^P+B\cdot5^P\le\mathrm{aTotal5}(P,t).
\end{equation*}
The cycle equation gives
\begin{equation*}
  \mathrm{aTotal5}(P,t)=5m(2^T-5^P),
\end{equation*}
hence
\begin{equation*}
  5^P+B\cdot5^P\le5m(2^T-5^P).
\end{equation*}
If the right-hand side is below $2\cdot5^P$, then $B=0$ and every
proper prefix satisfies $2^{\wordWeight_j}<5^j$.  These are
\texttt{pmi\_b\_count} and \texttt{pmi\_b\_no\_bad\_prefix}.

\subsection{Margin balance}

Let $C_1,C_2,C_3$ be the counts of steps of weight $1$, $2$, and at
least $3$.  The nonnegative-margin condition gives
\begin{equation*}
  (3-\alpha)C_3\le(\alpha-1)C_1+(\alpha-2)C_2,
\qquad \alpha=\log_2 5,
\end{equation*}
with the strict version used for the small-budget cycle words.

\subsection{Rise steps and residues modulo 5}

A rise step of weight $1$ satisfies
\begin{equation*}
  2y=5x+1,
\end{equation*}
so $y\equiv3\pmod5$.  A rise step of weight $2$ satisfies
\begin{equation*}
  4y=5x+1,
\end{equation*}
so $y\equiv4\pmod5$.  These rules are formalised as
\texttt{rise\_step\_next\_mod\_five} and feed the block-head residue
lemma of the cycle bridge.

\subsection{Derivation of the PMI identity}

Write the word molecule as
\begin{equation*}
  \wordA(w)=\sum_{j=0}^{L-1}2^{\wordWeight_j}\,5^{L-1-j}.
\end{equation*}
Divide $5\wordA(w)$ by $5^L$:
\begin{equation*}
  \frac{5\wordA(w)}{5^L}
    =\sum_{j=0}^{L-1}2^{\wordWeight_j}\,5^{-j}.
\end{equation*}
Since
\begin{equation*}
  2^{-m_j}=2^{\wordWeight_j-j\log_2 5}
    =2^{\wordWeight_j}\,5^{-j},
\end{equation*}
the identity follows.

The same derivation for a closed cycle word with cycle state $m$ and
period $P$ gives the algebraic form
\begin{equation*}
  \mathrm{aTotal5}(P,t)=5m(2^T-5^P),
\end{equation*}
where $T=\wordWeight(w)$ is the total cycle weight.

\subsection{Worked example}

The word $[2,1]$ is valid from $7$, with
\begin{equation*}
  7 \xrightarrow{t=2} 9 \xrightarrow{t=1} 23,
\end{equation*}
so $\wordOrbit([2,1],7)=23$.  Its word molecule is
\begin{equation*}
  \wordA([2,1])=2^0\cdot5^1+2^2\cdot5^0=9,
\end{equation*}
and
\begin{equation*}
  2^3\cdot23=184=5^2\cdot7+9.
\end{equation*}
The word $[1,2]$ is valid from $9$:
\begin{equation*}
  9 \xrightarrow{t=1} 23 \xrightarrow{t=2} 29,
\qquad
  \wordA([1,2])=2^0\cdot5^1+2^1\cdot5^0=7,
\end{equation*}
with $2^3\cdot29=232=5^2\cdot9+7$.

The word $[1,1]$ is valid from $9$ in the divisibility sense but is
not the exact accelerated word: the second weight at $23$ is
$v_2(116)=2$, so $\wordOrbit([1,1],9)=58$ is an intermediate state not
on the accelerated orbit of $7$.  This is the distinction between
$\GeneralOrbit$ and $\FullOrbit$.

\subsection{Worked PMI example}

For $w=[2,1,2]$ from $7$, we have $L=3$, $\wordWeight=5$, and
$\wordA(w)=53$, with
\begin{equation*}
  2^5\cdot29=928=5^3\cdot7+53.
\end{equation*}
The prefix weights are $\wordWeight_0=0$, $\wordWeight_1=2$,
$\wordWeight_2=3$, so
\begin{equation*}
  \sum_{j=0}^{2}2^{-m_j}
    =1+\frac45+\frac8{25}
    =\frac{53}{25}
    =\frac{5\wordA(w)}{5^3}.
\end{equation*}

\subsection{Valuation toolkit}

The following elementary facts are used throughout the proof:
\begin{itemize}[leftmargin=*]
\item $n=2^{\vtwo{n}}\operatorname{oddpart}(n)$ with odd part odd;
\item if $u$ is odd, then $\vtwo{2^k u}=k$;
\item for odd $x$, $\vtwo{5x+1}\ge1$;
\item $T(x)\le(5x+1)/2$;
\item a C3 step from $x\ge1$ satisfies $(5x+1)/2^t<x$.
\end{itemize}
The corresponding Lean names are listed in the main paper.

\subsection{Notation summary}

\begin{center}
\small
\begin{tabular}{@{}p{0.14\textwidth}p{0.80\textwidth}@{}}
\toprule
Symbol & Meaning \\
\midrule
$j$ & reset step depth \\
$k_0$ & exponent in $s5^{k_0}=r+1$ \\
$t$ & reset step weight, $t\in\{1,2\}$ \\
$\delta$ & reset parameter: $1$ for $t=1$, $\{1,3\}$ for $t=2$ \\
$s$ & odd terminal part with $5\nmid s$ \\
$e$ & incoming step weight, $e\ge2$ \\
$g$ & full-orbit state $\Orbit_7(j-1)$ \\
$x$ & candidate predecessor $2^{e-1}g+\delta5^{j-1}$ \\
$r_j$ & block head after the reset step \\
$r_s$ & terminal state of the block tail \\
$t_j$ & step weight at depth $j$ \\
$\wordWeight_j$ & prefix weight at depth $j$ \\
$L$ & block-tail length \\
$H_s$ & unified-core threshold parameter \\
\bottomrule
\end{tabular}
\end{center}

\subsection{Block equation}

For a block with step weights $u_1,\dots,u_L$ and total weight $U$,
\begin{equation*}
  2^U r_L=5^L r_0+B,
\end{equation*}
where $B$ is the block molecule.  The reset block satisfies
\begin{equation*}
  t=2:\quad 4(r_j+1)=5^{k_0+1}s+\delta5^j,
\end{equation*}
\begin{equation*}
  t=1:\quad 2(r_j+1)=5^{k_0+1}s+5^j-2.
\end{equation*}
These equations are used by the corrected window bound.  The block
equation is exactly the word orbit identity applied to the block word;
no averaging or approximation is involved.

For example, the block starting at $29$ with weights $[1,1,2]$ has
$L=3$, $U=4$, and $B=39$, with
\begin{equation*}
  2^4\cdot229=3664=5^3\cdot29+39.
\end{equation*}
